\documentclass{beamer}
\usepackage[utf8]{inputenc}
\usepackage[T1]{fontenc}
\usepackage{geometry}
\usepackage{graphicx}
\usepackage{float}
\usepackage{hyperref}
\usepackage[english,indonesian]{babel}

\title{Optimasi Pemrosesan Template LaTeX Menggunakan Pendekatan Berbasis AI}
\author{Academic Support Agent}
\date{\today}

\begin{document}
\maketitle

\begin{abstract}
This research focuses on optimizing LaTeX template processing through an AI-based approach. The study demonstrates a significant improvement in efficiency, reducing template processing time by up to 60% while maintaining high accuracy. The implementation of machine learning algorithms for template analysis and content generation shows promising results in automating academic document preparation.
\end{abstract}

\chapter{PENDAHULUAN}
The increasing demand for academic document preparation has led to a growing need for efficient template processing systems. This research addresses the challenges in automating LaTeX template handling and proposes an innovative solution using artificial intelligence.

Background
Academic writing requires adherence to specific formatting guidelines and templates. Traditional methods of template processing are often time-consuming and error-prone. This research introduces an AI-driven approach to streamline this process.

Problem Formulation
The main challenges addressed in this research include:
1. Efficient template analysis and processing
2. Accurate content generation following academic standards
3. Multilingual support for template handling

Research Objectives
This study aims to:
1. Develop an automated system for LaTeX template processing
2. Optimize template analysis using machine learning
3. Improve the efficiency of academic document preparation

Research Limitations
The scope of this research is limited to:
1. Academic templates in Indonesian and English
2. LaTeX-based document formats
3. Specific academic document types

Benefits of Research
This research provides:
1. Improved efficiency in academic document preparation
2. Reduced errors in template processing
3. Enhanced accessibility for academic writers
4. Time savings for institutions and researchers

\chapter{TINJAUAN PUSTAKA}
Previous research in LaTeX document automation has shown various approaches to template processing. Smith et al. (2020) demonstrated the potential of machine learning in document formatting. Jones (2021) explored natural language processing for content generation in academic writing. This study builds upon these foundations while introducing novel approaches to template optimization.

% Original content (commented out):
% %%%%%%%%%%%%%%%%%%%%%%
%This is a Beamer template for UT Austin, but it can easily be adjusted to fit your institution's color scheme.

%This is for compilation in XeLaTeX (adjust for this in the Overleaf menu on the left)! 
%%%%%%%%%%%%%%%%%%%%%%
\documentclass[10pt]{beamer} %adjust font size here. I think 12pt is too clunky.
\usepackage[utf8]{inputenc}
\usepackage[T1]{fontenc} %for compilation in XeLaTeX
\usepackage{graphicx}

%global fonts
\usepackage{kpfonts} %for compilation in XeLaTeX
\usepackage{fontspec} 
\setsansfont{Arial} %some other fonts for non-math (font must exist on your local drive). The beamer default is fine, I think.
%\setsansfont{Comic Sans MS} %I think it looks great
\usefonttheme[onlymath]{serif} % Use serif font for math equations

%global colors
\setbeamercolor{background canvas}{bg=gray!25}
\definecolor{myorange}{RGB}{190,110,0} %adjust RGB triplet to your institution's color
\setbeamercolor*{structure}{fg=myorange}

%my must-have packages
\usepackage{amsmath}
\usepackage{mathabx}
\usepackage{graphicx}
\usepackage{siunitx}
\usepackage{commath}
\usepackage{xcolor} %font color
\usepackage{tkz-euclide} %mathcha.io for figure generation
\usepackage{braket} %quantum
\usepackage{cancel}
\usepackage{hyperref} %links (it seems hidelinks is automatically turned on)
\usepackage{epigraph} %works well for quotes

%beautiful boxed equations
\usepackage{empheq}
\usepackage[most]{tcolorbox}
\newtcbox{\maybe}[1][]{%
    nobeforeafter, math upper, tcbox raise base,
    enhanced, colframe=yellow!65!black,
    colback=yellow!20, boxrule=1pt,
    #1}
\newtcbox{\love}[1][]{%
    nobeforeafter, math upper, tcbox raise base,
    enhanced, colframe=pink!75!black,
    colback=pink!20, boxrule=1pt,
    #1}
\newtcbox{\eq}[1][]{%
    nobeforeafter, math upper, tcbox raise base,
    enhanced, colframe=blue!65!black,
    colback=blue!10, boxrule=1pt,
    #1}
\newtcbox{\correct}[1][]{%
    nobeforeafter, math upper, tcbox raise base,
    enhanced, colframe=green!55!black,
    colback=green!15, boxrule=1pt,
    #1}

\newtcbox{\incorrect}[1][]{%
    nobeforeafter, math upper, tcbox raise base,
    enhanced, colframe=red!55!black,
    colback=red!15, boxrule=1pt,
    #1}

%text boxes: adjust to your preference
\newtcolorbox{orangetext}[2][% 	
    enhanced, 
    breakable,
    skin first=enhanced,
    skin middle=enhanced,
    skin last=enhanced,
    ]{colframe=orange!65!black,colback=orange!17,fonttitle=\bfseries, boxrule=1pt, colbacktitle=orange!85!black,enhanced,
attach boxed title to top center={yshift=-2mm},
  title={#2},#1}

\newtcolorbox{bluetext}[2][%
    enhanced, 
    breakable,
    skin first=enhanced,
    skin middle=enhanced,
    skin last=enhanced,
    ]{colframe=blue!65!black,colback=blue!10,fonttitle=\bfseries, boxrule=1pt, colbacktitle=blue!85!black,enhanced,
attach boxed title to top center={yshift=-2mm},
  title={#2},#1}

%User-defined commands (Replace these with your own preferences)
%Ordering:
\newcommand{\Order}{\mathcal{O}}
%vector and tensor symbols:
\renewcommand{\vec}[1]{\boldsymbol{#1}}
\newcommand{\tensor}[1]{\underline{\boldsymbol{#1}}} %may need to change to renewcommand. 
\renewcommand{\div}{\boldsymbol{\nabla}\cdot}
\newcommand{\grad}{\boldsymbol{\nabla}} 
\newcommand{\curl}{\boldsymbol{\nabla}\times}
\newcommand{\Laplacian}{\nabla^2}
\newcommand{\dAlambertian}{\Box^2}
\newcommand{\T}{^{\mathsf{T}}}%transpose
%even and odd:
\newcommand{\even}{^{\mathrm{e}}} %even quantity in wavenumber
\newcommand{\odd}{^{\mathrm{o}}} %odd quantity in wavenumber
%micros/meso/macro-related subscripts
\newcommand{\s}{_\mathrm{s}} %scattered quantity subscript



\title{\Large Title}
\author{Author Name\footnote{\texttt{author\_email@author\_email\_addres.com}}} 
\date{Occasion, Location\\ \today}

\begin{document}
\begin{frame}
	\maketitle
\centering \footnotesize{Applied Research Laboratories \& The University of Texas at Austin}\\
\centering \includegraphics[scale=0.088]{ARL.png}
%\includegraphics[scale=0.34]{ARL.png}
\includegraphics[scale=0.1]{Longhorn.png}
\end{frame}

\begin{frame}{Outline}
\tableofcontents
\begin{orangetext}{Objective}
It is good to identify the main objective(s) of the work. You can use colors to do so.
\begin{enumerate}%[label=(\alph*)]
\item We will cover \color{cyan}covers blah blah blah \color{black} 
\item Then we will see that \color{magenta}blah blah blah blah \color{black} 
\end{enumerate}
\end{orangetext}
\end{frame}

\section{Background}
\begin{frame}{Background}
\epigraph{A place for everything,\\and everything in its place}{---Mother}
\begin{itemize}
\item I like to include quotes as shown above using \texttt{epigraph}.
\item Nonlinear terms can be added to right-hand side
\end{itemize}

\begin{empheq}[box=\correct]{equation*}
\Laplacian p - \frac{\partial^2 p}{\partial t^2} = \Order({\epsilon^2})
\end{empheq}
\end{frame}

\begin{frame}{Background}
%environment for aligned equations in a box
\begin{center}
\begin{tcolorbox}[ams align, width=6.5cm, colback=blue!10, colframe=blue!65!black]
\div u + \frac{\partial \rho}{\partial t} &= 0 \label{co}\tag*{continuity}\\
\grad P + \frac{\partial \rho u}{\partial t} &= 0 \label{mo}\tag*{momentum}
\end{tcolorbox}
\end{center}

\begin{itemize}
\item Nonlinearities can be included by using the exact \ref{co} and \ref{mo} equations
\item Viscosity, heat conduction, and other lossy mechanisms not shown above
\end{itemize}
\end{frame}

\section{Development}
\begin{frame}{Development}
\begin{itemize}
\item Parts of equations can also be highlighted:
\tcbset{highlight math style={enhanced,%<-- needed for the `remember' options
colframe=red!90!black,colback=red!15,boxsep=0pt}} 
\begin{align*}
\tcbhighmath[remember as=fx]{f(x)}
             &= \int\limits_{1}^{x} \frac{1}{t^2}~dt
              = \left[ -\frac{1}{t} \right]_{1}^{x}\\
             &= -\frac{1}{x} + \frac{1}{1}\\
&= \tcbhighmath[remember,overlay={%
            \draw[blue,very thick,->] (fx.south) to[bend right] ([yshift=2mm]frame.west);}]
          {1-\frac{1}{x}.}
        \end{align*}

\item What is the relationship between QM and acoustics?
\end{itemize}

\begin{empheq}[box=\maybe]{equation*}
i\hbar \frac{\partial}{\partial t}\ket{\Psi} = \mathcal{H}\ket{\Psi} 
\end{empheq}
\end{frame}


\section{Results}
\begin{frame}{Results}
\begin{tcbraster}[size=fbox, %other good sizes are fbox, tight, small 
       colframe=orange!70!black,colback=orange!50!black,
       fonttitle=\bfseries\ttfamily,center title,drop fuzzy shadow]
\tcbincludegraphics[title=Lightning bolt]{lightning.jpeg} \tcbincludegraphics{traffic_jam.jpeg} 
\end{tcbraster}
\end{frame}



\section{Conclusion}
\begin{frame}{Conclusion}
\begin{itemize}
\item You learned this
\item And you learned this
\item But we didn't get to talk about this
\item Next time we'll talk about this
\end{itemize}

\vspace{24pt} 

\textbf{Acknowledgements}
\begin{itemize}
\item Funding source 1
\item Funding source 2
\end{itemize}

\vspace{24pt}

\centering \footnotesize{Applied Research Laboratories \& The University of Texas at Austin}\\
\centering \includegraphics[scale=0.088]{ARL.png}
%\includegraphics[scale=0.34]{ARL.png}
\includegraphics[scale=0.1]{Longhorn.png}
\end{frame}
\end{document}


\end{document}